\par {\bfseries ВВЕДЕНИЕ \dotfill 5}
\par {\bfseries 1 АНАЛИЗ ПРЕДМЕТНОЙ ОБЛАСТИ И ПОСТАНОВКА ЗАДАЧИ \dotfill 6}
\par 1.1 Общие сведения о сервисах заказа такси \dotfill 6
\par 1.2 Особенности разработки высоконагруженных серверных систем \dotfill 8
\par 1.3 Сравнительный анализ монолитной и микросервисной архитектур \dotfill 10
\par 1.4 Анализ существующих решений на рынке такси-сервисов \dotfill 12
\par 1.5 Требования к серверной части приложения \dotfill 15
\par 1.6 Постановка задачи \dotfill 17
\par {\bfseries 2 ПРОЕКТИРОВАНИЕ СИСТЕМЫ \dotfill 19}
\par 2.1 Проектирование общей архитектуры системы \dotfill 19
\par 2.2 Проектирование клиентской части \dotfill 20
\par 2.3 Архитектура взаимодействия компонентов системы \dotfill 22
\par 2.4 Проектирование структуры базы данных \dotfill 25
\par 2.5 Разработка API и протоколов обмена данными \dotfill 28
\par {\bfseries 3 РЕАЛИЗАЦИЯ ОСНОВНЫХ МОДУЛЕЙ СИСТЕМЫ \dotfill 32}
\par 3.1 Реализация сервиса управления заказами \dotfill 32
\par 3.2 Разработка модуля обработки геолокации и поиска ближайших водителей \dotfill 35
\par 3.3 Система авторизации и аутентификации на основе JWT \dotfill 38
\par 3.4 Интеграция брокера сообщений для обработки событий в реальном времени \dotfill 40
\par {\bfseries 4 ТЕСТИРОВАНИЕ И РАЗВЕРТЫВАНИЕ СИСТЕМЫ \dotfill 43}
\par 4.1 Модульное и интеграционное тестирование компонентов \dotfill 43
\par 4.2 Контейнеризация приложения с использованием Docker \dotfill 44
\par 4.3 Настройка CI/CD процессов и мониторинг серверной части \dotfill 46
\par {\bfseries 5 ТЕХНИКО-ЭКОНОМИЧЕСКОЕ ОБОСНОВАНИЕ \dotfill 48}
\par 5.1 Исходные данные \dotfill 48
\par 5.2 Расчет объёма функций программного обеспечения \dotfill 48
\par 5.3 Расчёт полной себестоимости программного модуля \dotfill 49
\par 5.4 Расчет цены и прибыли по программному продукту \dotfill 54
\par {\bfseries ЗАКЛЮЧЕНИЕ \dotfill 56}
\par {\bfseries СПИСОК СОКРАЩЕНИЙ \dotfill 57}
\par {\bfseries СПИСОК ИСПОЛЬЗОВАННЫХ ИСТОЧНИКОВ \dotfill 58}
\par {\bfseries ПРИЛОЖЕНИЕ А: ТЕКСТ ПРОГРАММЫ}
